% !PRINTED 664
\weberpage{664}
%<<CONT
\gesperrt{so we assign to such a coincidence of values
a point $\mathfrak{P}$.} We say: $\alpha = \alpha_0$, or $\alpha$ has the value $\alpha_0$
at the point $\mathfrak{P}$. Two points $\mathfrak{P}$ and $\mathfrak{P}'$ are called different then and only
then, when there is at least one function in $\Omega$ that has
different values at $\mathfrak{P}$ and $\mathfrak{P}'$.

Rules 2. to 5. fail in the exceptional cases, namely
2. and 3., when $\alpha_0$ and $\beta_0 = \infty$; 4., when $\alpha_0 = 0$, $\beta_0 = \infty$
or $\alpha_0 = \infty$, $\beta_0 = 0$; and 5., when $\alpha_0$ and $\beta_0$ are both $= 0$
or both $= \infty$.

Nevertheless the functions $(\alpha \pm \beta)_0$, $(\alpha\beta)_0$ and $(\alpha : \beta)_0$
must have determinate values in these cases too; these values, however, are not
 determined directly from prescriptions 2. to 5., but in an indirect
way.

\gesperrt{The point is accordingly an invariant concept belonging
 to the field $\Omega$}, one that has nothing to do with the accidental
circumstance of which of the functions of $\Omega$ we regard as the
independent variable.

To find all points, one proceeds as follows. If $\mathfrak{P}$ is a point,
then there are functions that have a finite value at it,
for if a function $\alpha$ is infinite at $\mathfrak{P}$, then $1 : \alpha$ is finite.
Such a function is to be taken as the independent variable $z$
and its value at the point $\mathfrak{P}$ denoted by $c$. All integral functions
 of $z$ are then likewise finite at $\mathfrak{P}$, for if $\omega$ is an
integral function that satisfies the equation
\[
\omega^{n} + a_1\omega^{n-1} + \cdots + a_n = 0
\]
then
\[
1 + \frac{a_1}{\omega} + \frac{a_2}{\omega^2} + \cdots + \frac{a_n}{\omega^n} = 0
\]
must be satisfied at $\mathfrak{P}$, and since the $a_1, a_2, \dots, a_n$, as integral rational
functions of $z$, are finite at $\mathfrak{P}$, $\omega$ cannot be $= \infty$.

Every congruence with respect to the modulus $z - c$ must therefore pass over at the
point $\mathfrak{P}$ into a correct equation, even when this
congruence subsists between holomorphic integral functionals.
If, then, $\tau$ is the basis form of $\Omega$ in $Z$, the
equation must hold, by § 184, (20):
\[
(\tau - \tau_1)(\tau - \tau_2)\dots(\tau - \tau_n) = 0.
\]
One of the factors of the left side must therefore vanish in $\mathfrak{P}$.
If this is $\tau - \tau_1$, then to this factor there corresponds a
prime factor $\pi$ of $z - c$. Thus in $\mathfrak{P}$ there will be

\weq{2}{\omega_1 = c_1, \quad \omega_2 = c_2, \dots, \omega_n = c_n,}
%CONT>>
